% Slide 06: Probabilistic Modeling: Sequence vs. Token
\begin{frame}{Probabilistic Modeling: Sequence vs.\ Token}
  \vspace{-0.1cm}
  {\footnotesize\textbf{1. RAG-Sequence (Sequence-Level Marginalization):}\par}
  \vspace{-4pt}
  \[
    p_{\mathrm{Seq}}(y \mid x) \approx \sum_{z \in \mathrm{top-}k} p_\eta(z \mid x) \prod_{i=1}^{T} p_\theta(y_i \mid x, z, y_{<i})
  \]
  \vspace{-4pt}
  {\footnotesize\textbf{2. RAG-Token (Token-Level Step Marginalization):}\par}
  \vspace{-4pt}
  \[
    p_{\mathrm{Tok}}(y \mid x) \approx \prod_{i=1}^{T} \left[ \sum_{z \in \mathrm{top-}k} p_\eta(z \mid x) p_\theta(y_i \mid x, z, y_{<i}) \right]
  \]

  \begin{itemize}
    \setlength{\itemsep}{0.12cm}
    \item \textbf{Latent Passage ($z$):} Retrieved documents act as latent variables marginalized across top-$k$ candidates.
    \item \textbf{Sequence Formulation:} Generates entire response $y$ from a single document $z$; preserves topic coherence.
    \item \textbf{Token Formulation:} Dynamically switches documents across tokens; ideal for synthesizing multi-source facts.
    \item \textbf{Variables:} $x$: query; $y_{<i}$: generated prefix; $\eta, \theta$: retriever/generator parameters; $p \in [0, 1]$.
  \end{itemize}

  \source{Lewis et al. (NeurIPS 2020), \S2.1, Equations (1)--(2).}
\end{frame}
